%%%%%%%%%%%%%%%%%%%%%%%%%%%%%%%%%%%%%%%%%%%%%%%%%%%%%%%%%%%%%%
%%%%%%%% MA-0801 Comunicación en las ciencias %%%%%%%%%%%%%%%%
%%%%%%%%%%%%%%%%%%%%%%%%%%%%%%%%%%%%%%%%%%%%%%%%%%%%%%%%%%%%%%
\newpage
\subsection*{MA-0801 Comunicación en las ciencias}
\addcontentsline{toc}{subsection}{MA-0801 Comunicación en las ciencias}

\hypertarget{MA0801}{}
\begin{refsection}
\begin{table}[H]
\centering{
\begin{tabular}{|ll|}
\hline
\textbf{Año:} IV  & \textbf{Requisitos:} Haber aprobado al menos uno de: MA-0661 \\ & Algebra Abstracta II, MA-0635 Introducción a la Topología \\ & o MA-0641 Matemática Computacional II \\ 
\textbf{Ciclo:} VIII  & \textbf{Correquisitos:} Ninguno \\ 
\textbf{Tipo de curso:} Seminario & \textbf{Horas presenciales por semana:} 2 \\ 
\textbf{Créditos:} 2 & \textbf{Horas de trabajo independiente por semana:} 4 \\ \hline
\end{tabular}
}
\end{table}

\subsubsection*{Descripción del curso}

Este curso está dirigido a personas estudiantes de los últimos ciclos de la carrera Bachillerato en Matemática. Su propósito es desarrollar habilidades de \textbf{escritura} y \textbf{expresión oral} para comunicar conocimiento científico --en particular matemático-- en una amplia variedad de contextos profesionales y académicos: docencia, propuestas de investigación, artículos y informes técnicos, solicitudes de financiamiento, presentaciones ante públicos especializados o generalistas, y otros formatos propios de la práctica matemática contemporánea.

El curso es \textbf{eminentemente práctico y participativo}: la persona estudiante prepara textos y presentaciones, las expone en clase y recibe retroalimentación de la persona docente y del grupo. Complementariamente, se incorporan lecturas y material audiovisual para trabajo extraclase. Un componente transversal es la \textbf{ética profesional}, incluida la formación en conducta responsable en la investigación (\emph{Responsible Conduct of Research}, RCR): integridad académica, atribución correcta de fuentes, manejo de datos y comunicación honesta de resultados.

El diseño del curso se alinea con competencias de comunicación científica reconocidas internacionalmente (redacción técnica, exposición oral, composición tipográfica con \LaTeX{}, búsqueda bibliográfica y planificación de carrera), adaptadas al contexto del Departamento de Matemáticas y, en perspectiva, a iniciativas de formación comunes en la Facultad de Ciencias.

\subsubsection*{Objetivos}

\textbf{General:} Que la persona estudiante sea capaz de comunicarse en forma oral y escrita para transmitir de forma efectiva ideas, resultados y argumentos matemáticos con claridad, rigor y pertinencia en contextos académicos y profesionales.

\textbf{Específicos:} Al finalizar el curso se espera que la persona estudiante sea capaz de:

\begin{enumerate}
\item Redactar informes y textos matemáticos técnicos con estructura clara, notación consistente y argumentación comprensible para el público previsto.
\item Planificar, preparar y exponer oralmente resultados o temas matemáticos, organizando apoyos visuales y respondiendo preguntas del auditorio.
%\item Componer documentos matemáticos con el sistema tipográfico \LaTeX{} u otra herramienta equivalente aprobada por la persona docente.
\item Localizar literatura pertinente, citarla y referenciarla conforme a normas académicas y principios de integridad intelectual.
\item Reconocer y aplicar criterios básicos de conducta responsable en la investigación y la comunicación científica.
\item Elaborar materiales orientados a la empleabilidad matemática (por ejemplo, currículum vitae y carta de presentación ante convocatorias con perfil cuantitativo).
\item Participar activamente en discusiones y retroalimentación entre pares sobre la calidad de exposiciones y textos científicos.
\end{enumerate}

\subsubsection*{Contenidos}

\begin{enumerate}

\item \textbf{Comunicación científica y ética profesional (2 semanas):}
\begin{enumerate}
    \item Roles de la comunicación en la práctica matemática e investigativa.
    \item Públicos, registros y propósitos comunicativos (especializado, general, docente, institucional).
    \item Integridad académica, plagio, autoría y conducta responsable en la investigación (RCR).
    \item Normas institucionales y buenas prácticas en el uso de datos, software y colaboración.
\end{enumerate}

\item \textbf{Escritura matemática técnica (3 semanas):}
\begin{enumerate}
    \item Principios de redacción clara: definiciones, teoremas, pruebas, ejemplos y contraejemplos.
    \item Estructura de informes, resúmenes (\emph{abstracts}) y notas de investigación.
    \item Propuestas de proyecto, solicitudes de apoyo y comunicación con comités evaluadores.
    \item Revisión por pares, retroalimentación y reescritura.
\end{enumerate}

%\item \textbf{Composición tipográfica con \LaTeX{} (2 semanas):}
%\begin{enumerate}
%    \item Entorno de trabajo, estructura de documentos y paquetes básicos.
%    \item Matemática display e inline, referencias cruzadas, bibliografía automatizada.
%    \item Figuras, tablas y buenas prácticas de legibilidad en documentos técnicos.
%\end{enumerate}

\item \textbf{Presentación oral y apoyos visuales (3 semanas):}
\begin{enumerate}
    \item Planificación de charlas cortas y presentaciones extensas; manejo del tiempo.
    \item Diseño de diapositivas o pizarra: jerarquía visual, ejemplos pertinentes, legibilidad.
    \item Exposición ante grupos pequeños y audiencias amplias; escucha activa y respuesta a preguntas.
    \item Retroalimentación constructiva entre pares.
\end{enumerate}

\item \textbf{Búsqueda bibliográfica y gestión de fuentes 2 semanas):}
\begin{enumerate}
    \item Bases de datos y recursos matemáticos (MathSciNet, zbMATH, arXiv, repositorios institucionales).
    \item Lectura crítica de artículos; síntesis y referencia adecuada.
    \item Herramientas básicas de gestión de referencias.
\end{enumerate}

\item \textbf{Comunicación en docencia y divulgación (2 semanas):}
\begin{enumerate}
    \item Explicación de ideas matemáticas a distintos niveles de preparación.
    \item Divulgación responsable: precisión sin sacrificar accesibilidad.
    \item Sensibilidad intercultural y trabajo con audiencias diversas en contextos científicos.
\end{enumerate}

\item \textbf{Planificación de carrera y comunicación profesional (2 semana):}
\begin{enumerate}
    \item Currículum vitae académico y perfil profesional en matemática.
    \item Cartas de presentación y comunicación escrita con empleadores o programas de posgrado.
    \item Entrevistas breves y \emph{elevator pitch} sobre intereses matemáticos.
\end{enumerate}

\item \textbf{Presentaciones integradoras del estudiantado (1 semana):}
\begin{enumerate}
    \item Exposiciones finales con retroalimentación formal (a través de rúbricas).
    \item Entrega de portafolio escrito (informe técnico o propuesta en \LaTeX{}).
\end{enumerate}

\end{enumerate}

\subsubsection*{Metodología}

El curso adopta la modalidad \textbf{seminario} con enfoque \textbf{práctico-participativo}. Las sesiones presenciales (2 horas semanales) combinan breves orientaciones de la persona docente con trabajo activo del estudiantado.

\noindent\textbf{Lineamientos metodológicos:} La persona docente a cargo del curso debe:

\begin{enumerate}
    \item Diseñar actividades en las que predomine la \textbf{práctica en clase}: presentaciones orales, discusiones guiadas, talleres de escritura y revisión entre pares.
    \item Asignar \textbf{lecturas y videos} para trabajo extraclase que preparen o consoliden cada bloque temático.
    \item Proporcionar \textbf{retroalimentación oportuna} sobre borradores escritos y exposiciones orales, con criterios explícitos de calidad.
    \item Integrar de forma transversal los temas de \textbf{ética e integridad académica} (RCR), vinculándolos a las tareas concretas del curso.
    \item Fomentar el respeto, la escucha activa y la crítica fundamentada en las discusiones grupales.
    \item \textbf{No utilizar exámenes} como instrumento principal de evaluación; la calificación debe basarse en productos comunicativos y participación sostenida.
\end{enumerate}

\textbf{Sugerencias metodológicas:}

\begin{enumerate}
    \item Alternar presentaciones individuales y en parejas o equipos pequeños.
    \item Incluir al menos una actividad de simulación (defensa oral breve, \emph{mock} de propuesta o revisión de un artículo).
    \item Invitar, cuando sea posible, a un profesional o investigador invitado para comentar sobre comunicación matemática en la práctica.
    \item Utilizar rúbricas compartidas con el estudiantado para autoevaluación y coevaluación.
    \item Recomendar el uso de software libre para edición tipográfica y gestión bibliográfica.
\end{enumerate}

% \subsubsection*{Evaluación}

% Se sugiere una evaluación \textbf{continua} basada en productos y participación, sin exámenes tradicionales:

% \begin{enumerate}
% \item \textbf{Presentaciones orales} (aprox.\ 35\%): al menos dos exposiciones durante el semestre, incluida una integradora final.
% \item \textbf{Trabajos escritos} (aprox.\ 40\%): informe técnico, propuesta o artículo breve en \LaTeX{}, con revisiones sucesivas.
% \item \textbf{Participación en clase} (aprox.\ 25\%): asistencia, discusiones, retroalimentación a pares y cumplimiento de lecturas asignadas.
% \end{enumerate}

% Los porcentajes son orientativos; la persona docente puede ajustarlos al inicio del curso, manteniendo la evaluación formativa y práctica como eje central.

\subsubsection*{Bibliografía}

\textbf{Referencias generales:}

\begin{enumerate}
\item Higham, N. J. (1998). \emph{Handbook of Writing for the Mathematical Sciences} (2nd ed.). SIAM.
\item Knuth, D. E., Larrabee, T., y Roberts, P. M. (1989). \emph{Mathematical Writing}. Mathematical Association of America.
\item Lamport, L. (1994). \emph{\LaTeX: A Document Preparation System} (2nd ed.). Addison-Wesley.
\item Council of Graduate Schools / responsables nacionales de política de investigación. Materiales sobre conducta responsable en la investigación (RCR).
\item Recursos en línea de la \emph{American Mathematical Society} y la \emph{Mathematical Association of America} sobre exposición y publicación matemática.
\end{enumerate}

\textbf{Lecturas y material audiovisual:} la persona docente seleccionará, por ciclo lectivo, artículos modelo, guías institucionales y videos sobre comunicación científica, los cuales serán anunciados al inicio del curso.

\end{refsection}
